\chapter{Analiza wymagań systemowych}

\section{Wymagania funkcjonalne}
System „French Riviera” został zaprojektowany jako wewnętrzna aplikacja hotelowa wspierająca obsługę gości, realizację usług oraz organizację pracy personelu. Zakres projektu skupia się na przygotowaniu widoków dla poszczególnych grup użytkowników, z uwzględnieniem ich codziennych zadań oraz najważniejszych informacji prezentowanych w interfejsie.

\subsection{Moduł Gościa Hotelowego (Tablet pokojowy)}
Interfejs przeznaczony dla gościa hotelowego, umożliwiający wygodne korzystanie z usług hotelowych za pośrednictwem tabletu znajdującego się w pokoju.
\begin{itemize}
    \item \textbf{Room Service:} Przeglądanie menu, wybór produktów, określanie ilości oraz dodawanie uwag do zamówienia.
    \item \textbf{Zgłoszenia pokojowe:} Szybkie zgłaszanie potrzeby sprzątania, wymiany ręczników, wymiany pościeli, uzupełnienia minibaru oraz usterek technicznych.
    \item \textbf{Usługi dodatkowe:} Przeglądanie i wybór dodatkowych atrakcji hotelowych, takich jak kort tenisowy, trening na siłowni, zabiegi SPA lub inne usługi dostępne w hotelu.
    \item \textbf{Panel informacyjny:} Podgląd podstawowych informacji o pobycie, takich jak numer pokoju, data wymeldowania, aktywne usługi i bieżące zgłoszenia.
    \item \textbf{Statusy:} Podgląd aktualnego statusu zamówień, zgłoszeń i wykupionych usług.
\end{itemize}

\subsection{Moduł Housekeepingu (Aplikacja mobilna)}
Interfejs mobilny przeznaczony dla obsługi pięter, skupiony na szybkiej realizacji zadań związanych z pokojami.
\begin{itemize}
    \item \textbf{Lista zadań:} Widok pokoi i zadań przypisanych do obsługi, z możliwością podziału według piętra, priorytetu lub rodzaju zadania.
    \item \textbf{Typy zadań:} Obsługa zadań takich jak sprzątanie pokoju, wymiana ręczników, wymiana pościeli oraz uzupełnienie minibaru.
    \item \textbf{Szczegóły zadania:} Podgląd numeru pokoju, rodzaju zgłoszenia, czasu zgłoszenia, priorytetu oraz dodatkowych uwag.
    \item \textbf{Aktualizacja statusu:} Oznaczanie zadania jako będącego w realizacji lub wykonanego.
\end{itemize}

\newpage

\subsection{Moduł Konserwatora / Maintenance (Aplikacja mobilna)}
Interfejs mobilny przeznaczony dla pionu technicznego, służący do obsługi usterek zgłaszanych w hotelu.
\begin{itemize}
    \item \textbf{Lista usterek:} Widok zgłoszeń technicznych przypisanych do pokoi.
    \item \textbf{Szczegóły zgłoszenia:} Podgląd numeru pokoju, rodzaju usterki, czasu zgłoszenia, priorytetu oraz dodatkowych uwag.
    \item \textbf{Aktualizacja statusu:} Oznaczanie zgłoszenia jako przyjętego do realizacji, będącego w naprawie lub zakończonego.
    \item \textbf{Filtrowanie zadań:} Przegląd zgłoszeń według statusu, priorytetu.
\end{itemize}

\subsection{Moduł Recepcji (Aplikacja desktopowa)}
Główny panel operacyjny służący do bieżącej obsługi pokoi i gości.
\begin{itemize}
    \item \textbf{Podgląd pokoi:} Widok pokoi z informacją o ich aktualnym statusie, np. wolny, zajęty, do sprzątania lub awaria.
    \item \textbf{Szczegóły pokoju:} Podgląd podstawowych informacji o pokoju, pobycie, aktywnych zgłoszeniach i usługach.
    \item \textbf{Check-in / Check-out:} Widoki wspierające obsługę meldunku i wymeldowania gościa.
\end{itemize}

\subsection{Moduł Kuchni (Tablet)}
Interfejs tabletowy przeznaczony do obsługi zamówień gastronomicznych składanych przez gości hotelowych.
\begin{itemize}
    \item \textbf{Kolejka zamówień:} Lista zamówień wraz z numerem pokoju, czasem złożenia i bieżącym statusem.
    \item \textbf{Szczegóły zamówienia:} Podgląd pozycji zamówienia, liczby produktów oraz dodatkowych uwag.
    \item \textbf{Aktualizacja statusu:} Oznaczanie zamówienia jako przygotowywanego lub gotowego.
\end{itemize}

\subsection{Moduł Menadżera (Aplikacja desktopowa)}
Panel zarządczy przeznaczony do organizacji pracy personelu i nadzoru nad realizacją zadań.
\begin{itemize}
    \item \textbf{Zarządzanie personelem:} Podgląd listy pracowników, podstawowych danych oraz organizacji pracy personelu.
    \item \textbf{Harmonogramy:} Tworzenie i edycja grafików zmianowych dla poszczególnych działów hotelu.
    \item \textbf{Przydzielanie zadań:} Delegowanie zadań do odpowiednich działów, takich jak housekeeping, maintenance czy kuchnia.
    \item \textbf{Podgląd realizacji:} Monitoring aktualnych zadań i opóźnień w ich wykonaniu.
\end{itemize}

\subsection{Moduł Administratora (Aplikacja desktopowa)}
Panel techniczny służący do zarządzania użytkownikami i utrzymania poprawnego działania systemu.
\begin{itemize}
    \item \textbf{Zarządzanie użytkownikami:} Tworzenie kont, edycja użytkowników, resetowanie haseł oraz nadawanie ról i uprawnień.
\end{itemize}

\section{Wymagania niefunkcjonalne}
Aby system był w pełni użyteczny w środowisku hotelowym, musi spełniać poniższe standardy projektowe:
\begin{itemize}
    \item \textbf{Responsywność:} Pełna adaptacja interfejsu do urządzeń mobilnych (Maintenance, Housekeeping) oraz tabletów i komputerów.
    \item \textbf{Intuicyjność:} Struktura nawigacji ograniczająca wysiłek poznawczy, co jest kluczowe przy dużej dynamice pracy personelu.
\end{itemize}